%! Author = sriadityavedantam
%! Date = 3/25/26

\section{Cantor's Theorem}

\lesson{5}{Friday 22 August 2025 9:10}{Day 5}

\begin{theorem}[Cantor's Diagonalization Argument]{
    Interval $C\subset\rr$ is uncountable.
}
	Suppose for the sake of contradiction, $C$ is countable, and let $c_1,c_2,\ldots$ enumerate $C$ such that each $c_i$ is written in decimal expansion:
    \begin{align*}
        c_1 &= (c_{11},c_{12},c_{13},\ldots)\\
        c_2 &= (c_{21},c_{22},\ldots)\\
        c_3 &= (c_{31},\ldots)\\
        &\vdots
    \end{align*}
    Construct $\gamma\in\rr$ such that $\gamma\neq c_i$ for all $i$
    \[
        \gamma_j \coloneqq
        \begin{cases}
            1 &, c_{jj} \neq 1\\
            0 &, c_{jj} = 1
        \end{cases}
    \]
    Then $\gamma$ differs from each $c_k$ in the $k$-th digit, so $\gamma\neq c_k$ for all $k$.
    This contradicts that $\{c_i\}$ enumerates $C$.
    Hence $C$ is uncountable.
\end{theorem}